% Apresentação Beamer da configuração 40.
% Estilo Madrid baseado em presentation_1.tex, fornecido pelo usuário.
% O conteúdo e os trechos C correspondem à versão atual do projeto.
% Compilar: pdflatex -interaction=nonstopmode -halt-on-error apresentacao.tex
\documentclass[aspectratio=169,11pt]{beamer}
\usetheme{Madrid}
\usepackage[T1]{fontenc}
\usepackage[utf8]{inputenc}
\usepackage[portuguese]{babel}
\usepackage{lmodern}
\usepackage{booktabs,array,listings}
\setbeamertemplate{navigation symbols}{}
\setbeamertemplate{footline}{\leavevmode\hbox{%
\begin{beamercolorbox}[wd=.77\paperwidth,ht=2.6ex,dp=1.2ex,leftskip=1em]{author in head/foot}%
Enthony | Douglas \hspace{1em} Tabela hash -- configuração 40
\end{beamercolorbox}%
\begin{beamercolorbox}[wd=.23\paperwidth,ht=2.6ex,dp=1.2ex,rightskip=1em plus1fil]{date in head/foot}%
\hfill\insertframenumber/\inserttotalframenumber
\end{beamercolorbox}}}
\setbeamersize{text margin left=9mm,text margin right=9mm}
\setlength{\emergencystretch}{2em}
\lstset{language=C,basicstyle=\ttfamily\small,keywordstyle=\color{blue!65!black},
commentstyle=\color{green!35!black},stringstyle=\color{red!55!black},
showstringspaces=false,columns=fullflexible,keepspaces=true,breaklines=true,
frame=single,rulecolor=\color{blue!25},aboveskip=5pt,belowskip=5pt,
literate={á}{{\'a}}1 {ã}{{\~a}}1 {â}{{\^a}}1 {à}{{\`a}}1
{é}{{\'e}}1 {ê}{{\^e}}1 {í}{{\'i}}1 {ó}{{\'o}}1 {õ}{{\~o}}1
{ô}{{\^o}}1 {ú}{{\'u}}1 {ç}{{\c c}}1 {Á}{{\'A}}1 {É}{{\'E}}1}
\newcommand{\arquivo}[1]{{\small\texttt{\detokenize{#1}}}}
\title[Tabela hash]{Tabela hash: índice de palavras}
\subtitle{Configuração 40: E1.1, H4, C2 e A4}
\author{Enthony Araujo de Oliveira \\ Douglas Patrick Cardozo Freitas}
\institute{Estrutura de Dados II\\Professor: Kennedy Reurison Lopes}
\date{}
\begin{document}
\begin{frame}\titlepage\end{frame}
\begin{frame}{Objetivo e roteiro}
\begin{block}{Objetivo do trabalho}
Implementar uma tabela hash própria e usá-la para consultar palavras de um texto.
\end{block}
\begin{enumerate}
\item Configuração e organização do projeto.
\item Normalização, conversão, função hash e colisões.
\item Operações e decisões da implementação.
\item Testes de correção e análise experimental.
\item Demonstração no terminal, com entradas e resultados esperados.
\end{enumerate}
\end{frame}

\begin{frame}{Nossa configuração: nº 40}
\renewcommand{\arraystretch}{1.45}
\begin{tabular}{@{}llp{8.7cm}@{}}\toprule
Código & Escolha & Significado no projeto\\\midrule
E1.1 & Vetor fixo & Capacidade escolhida ao criar, sem crescimento.\\
H4 & Dobra & Soma de grupos decimais de quatro dígitos.\\
C2 & Sondagem linear & Em colisão, procura a próxima posição.\\
A4 & Índice de palavras & Palavra normalizada como chave, frequência e linhas como valor.\\\bottomrule
\end{tabular}
\medskip
\begin{block}{Implementação própria em C11}
Usamos vetores, strings, leitura de arquivo e temporização. Não usamos uma tabela hash pronta.
\end{block}
\end{frame}

\begin{frame}{O que a aplicação A4 guarda}
\begin{columns}[T]
\column{.46\textwidth}
\textbf{Exemplo de texto}\par\medskip
\texttt{Linha 1: Hash hash!}\par
\texttt{Linha 2: HASH e tabela.}\par\medskip
A mesma palavra pode aparecer várias vezes, inclusive com caixa diferente.
\column{.49\textwidth}
\textbf{Depois da normalização}\par\medskip
\renewcommand{\arraystretch}{1.35}
\begin{tabular}{lcl}\toprule
Palavra & Frequência & Linhas\\\midrule
hash & 3 & 1, 2\\
e & 1 & 2\\
tabela & 1 & 2\\\bottomrule
\end{tabular}
\end{columns}
\bigskip
\begin{block}{Frequência e linhas são informações diferentes}
A frequência conta cada ocorrência. O vetor de linhas registra cada linha uma única vez durante a leitura sequencial.
\end{block}
\end{frame}

\begin{frame}{Organização do código}
\renewcommand{\arraystretch}{1.35}
\begin{tabular}{@{}lp{9.5cm}@{}}\toprule
Arquivo & Responsabilidade\\\midrule
\arquivo{src/tabela.h} & Tipos e declarações das operações.\\
\arquivo{src/tabela.c} & Conversão, dobra, vetor, sondagem, operações e exibição.\\
\arquivo{src/aplicacao.c} & Normalização e indexação do arquivo de texto.\\
\arquivo{src/main.c} & Menu, entrada do usuário e abertura do arquivo.\\
\arquivo{src/experimentos.c} & Medições experimentais. Correção pelo menu.\\\bottomrule
\end{tabular}\par
\medskip
\textbf{Fluxo:} o menu chama a aplicação, que normaliza as palavras e usa as operações da tabela.
\par\medskip
Os cabeçalhos permitem que os arquivos compartilhem os mesmos tipos e funções.
\end{frame}

\begin{frame}[fragile]{Estruturas usadas na tabela}
\begin{lstlisting}[language=C,basicstyle=\ttfamily\small]
typedef struct {
    char chave[TAM_PALAVRA];
    int frequencia;
    int *linhas;
    size_t quantidade_linhas;
} Palavra;

typedef struct {
    Estado estado;
    Palavra palavra;
} Posicao;
\end{lstlisting}
\begin{itemize}
\item \texttt{chave}: palavra de até 255 bytes, mais o terminador da string.
\item \texttt{linhas}: endereço do vetor de inteiros com as linhas do texto.
\item Cada posição contém um registro e um estado: vazio, ocupado ou removido.
\end{itemize}
\end{frame}

\begin{frame}{Normalização das palavras}
\begin{itemize}
\item Convertemos letras para minúsculas.
\item Mapeamos os acentos portugueses tratados no código para a letra base.
\item Na leitura do arquivo, espaços e pontuação separam as palavras.
\item O menu espera uma palavra por vez e aplica a mesma normalização.
\end{itemize}
\medskip
\renewcommand{\arraystretch}{1.3}
\begin{tabular}{ll}\toprule
Entrada & Chave usada\\\midrule
\texttt{Hash}, \texttt{HASH} & \texttt{hash}\\
\texttt{Colisão,} & \texttt{colisao}\\
\texttt{guarda-chuva} no arquivo & \texttt{guarda} e \texttt{chuva}\\\bottomrule
\end{tabular}
\medskip
\begin{alertblock}{Limites da regra}
Tokens com mais de 255 bytes são rejeitados. O tratamento de acentos é simples, sem suporte completo a todas as formas Unicode.
\end{alertblock}
\end{frame}

\begin{frame}[fragile]{Da palavra para um inteiro}
\begin{lstlisting}[language=C,basicstyle=\ttfamily\footnotesize]
static uint64_t chave_para_inteiro(const char *palavra)
{
    uint64_t numero = 0;
    for (size_t indice = 0; palavra[indice] != '\0'; indice++) {
        unsigned char byte = (unsigned char)palavra[indice];
        numero = numero * 131 + byte;
    }
    return numero;
}
\end{lstlisting}
\begin{itemize}
\item Cada byte influencia o resultado, sempre na mesma ordem.
\item Exemplo: \texttt{hash} produz \textbf{235481250}.
\item \texttt{uint64\_t} usa 64 bits sem sinal, com aritmética módulo $2^{64}$.
\item Essa conversão prepara a chave textual. A função atribuída continua sendo a dobra H4.
\end{itemize}
\end{frame}

\begin{frame}[fragile]{H4: cálculo por dobra}
\begin{lstlisting}[language=C,basicstyle=\ttfamily\footnotesize]
size_t hash_indice(const char *palavra, size_t capacidade)
{
    uint64_t numero = chave_para_inteiro(palavra);
    uint64_t soma = 0;
    while (numero > 0) {
        uint64_t grupo = numero % 10000;
        soma = soma + grupo;
        numero = numero / 10000;
    }
    return (size_t)(soma % capacidade);
}
\end{lstlisting}
Para \texttt{hash}: $235481250$ se divide em $2$, $3548$ e $1250$.
\[
\text{soma}=2+3548+1250=4800,\qquad h(\texttt{hash})=4800\bmod m.
\]
Com $m=2$, o índice é $0$. O módulo garante um resultado em $[0,m-1]$.
\end{frame}

\begin{frame}{E1.1: capacidade fixa e estados}
\begin{itemize}
\item Capacidade padrão: $m=64$. Podemos escolher outro valor antes de iniciar.
\item \texttt{calloc} reserva o vetor inicial. O programa não o redimensiona.
\item ``Estático'', na convenção da atividade, significa capacidade fixa durante o uso.
\end{itemize}
\medskip
\renewcommand{\arraystretch}{1.4}
\begin{tabular}{lp{10cm}}\toprule
Estado & Como interpretar\\\midrule
\texttt{VAZIO} & Nunca ocupado desde a criação. Uma busca pode parar aqui.\\
\texttt{OCUPADO} & Há uma palavra válida para comparar.\\
\texttt{REMOVIDO} & Um registro foi excluído. A busca deve continuar.\\\bottomrule
\end{tabular}\par
\medskip
O vetor de linhas de uma palavra pode crescer. Isso não aumenta as posições da tabela hash.
\end{frame}

\begin{frame}{C2: colisão e sondagem linear}
\[
p_i(k)=(h(k)+i)\bmod m,\qquad i=0,1,\ldots,m-1.
\]
Com $m=2$, \texttt{hash} e \texttt{tabela} têm índice inicial $0$.
\medskip
\begin{center}
\renewcommand{\arraystretch}{1.4}
\begin{tabular}{cccl}\toprule
Palavra & Índice inicial & Posição ocupada & Motivo\\\midrule
\texttt{hash} & 0 & 0 & Primeira posição livre.\\
\texttt{tabela} & 0 & 1 & A posição 0 já estava ocupada.\\\bottomrule
\end{tabular}
\end{center}
\begin{block}{O módulo permite voltar ao início}
Se a posição inicial for a última do vetor, a próxima tentativa usa a posição 0. Após m tentativas, todas as posições foram visitadas.
\end{block}
\end{frame}

\begin{frame}{Inserção: decisões do algoritmo}
\begin{enumerate}
\item Procurar a chave para impedir uma duplicata, inclusive depois de um removido.
\item Calcular o índice inicial com H4.
\item Percorrer o vetor pela sondagem linear.
\item Usar a primeira posição vazia ou removida.
\item Registrar palavra e frequência, incrementar n e contar a colisão aplicável.
\end{enumerate}
\begin{block}{Dois resultados importantes}
\textbf{Duplicata:} retorna \texttt{RES\_DUPLICADA}, sem substituir o valor.
\par\smallskip
\textbf{Sem posição disponível:} retorna \texttt{RES\_CHEIA}, sem crescer o vetor.
\end{block}
\end{frame}

\begin{frame}[fragile]{Busca: como o código percorre o vetor}
\begin{lstlisting}[language=C,basicstyle=\ttfamily\footnotesize]
size_t indice = (indice_inicial + tentativa) % tabela->capacidade;
Posicao *posicao = &tabela->posicoes[indice];
if (posicao->estado == VAZIO) {
    break;
}
if (posicao->estado == OCUPADO) {
    if (strcmp(posicao->palavra.chave, chave) == 0) {
        return indice;
    }
}
\end{lstlisting}
Trecho dentro do laço de sondagem. \texttt{strcmp(...)=0} significa strings iguais. Se não encontrar a chave, a função retorna \texttt{tabela->capacidade}. O laço visita no máximo m posições.
\end{frame}

\begin{frame}[fragile]{Remoção: por que manter REMOVIDO}
\begin{lstlisting}[language=C,basicstyle=\ttfamily\footnotesize]
size_t indice = localizar(tabela, chave);
if (indice == tabela->capacidade) {
    return RES_AUSENTE;
}
free(tabela->posicoes[indice].palavra.linhas);
tabela->posicoes[indice].palavra = (Palavra){0};
tabela->posicoes[indice].estado = REMOVIDO;
tabela->quantidade--;
return RES_OK;
\end{lstlisting}
\begin{itemize}
\item Após remover \texttt{hash} da posição 0, \texttt{tabela} ainda está na posição 1.
\item Se a posição 0 virasse \texttt{VAZIO}, a busca pararia antes de encontrar \texttt{tabela}.
\item \texttt{REMOVIDO} preserva o caminho e permite reutilizar a posição depois.
\end{itemize}
\end{frame}

\begin{frame}[fragile]{Atualização e listagem}
\begin{lstlisting}[language=C,basicstyle=\ttfamily\footnotesize]
size_t indice = localizar(tabela, chave);
if (indice == tabela->capacidade) {
    return RES_AUSENTE;
}
tabela->posicoes[indice].palavra.frequencia = frequencia;
return RES_OK;
\end{lstlisting}
\begin{itemize}
\item Exemplo: atualizar \texttt{tabela} de frequência 2 para 7 mantém a mesma chave.
\item A atualização mantém as linhas já registradas e não cria outro registro.
\item A listagem percorre o vetor e imprime somente posições ocupadas.
\item A exibição da estrutura também mostra vazios, removidos, índice inicial e contadores.
\end{itemize}
\end{frame}

\begin{frame}[fragile]{Indexação: palavra nova ou repetida}
\begin{lstlisting}[language=C,basicstyle=\ttfamily\small]
if (palavra != NULL) {
    palavra->frequencia++;
    tabela_adicionar_linha(palavra, numero_linha);
    ocorrencias_aceitas++;
}
\end{lstlisting}
\begin{itemize}
\item Primeiro, buscamos a palavra normalizada na tabela.
\item Se ainda não existe, tentamos inseri-la com frequência zero.
\item Com um registro disponível, somamos a ocorrência e guardamos a linha.
\item Palavra já existente recebe novas ocorrências mesmo com a tabela cheia.
\end{itemize}
\end{frame}

\begin{frame}{Contadores, capacidade e custo}
\begin{itemize}
\item $m$: quantidade fixa de posições do vetor.
\item $n$: quantidade atual de palavras distintas armazenadas.
\item $\alpha=n/m$: fator de carga, isto é, a ocupação da tabela.
\item Uma colisão é contada quando uma nova chave aceita encontra sua posição inicial ocupada por outra.
\item Duplicatas e atualizações não aumentam esse contador. Remover reduz n, mas não apaga o histórico de colisões.
\end{itemize}
\begin{block}{Custo das operações}
A sondagem pode visitar até m posições: pior caso $O(m)$, com chaves de tamanho limitado. A distribuição das chaves e a ocupação influenciam o custo observado.
\end{block}
\end{frame}

\begin{frame}{Testes manuais: entradas e resultados}
\small
\renewcommand{\arraystretch}{1.25}
\begin{tabular}{@{}p{6.6cm}p{6.5cm}@{}}\toprule
Entrada / caso & Esperado e obtido na execução\\\midrule
Buscar em tabela vazia & Ausente. OK.\\
Inserir \texttt{hash=1} e buscar & Frequência 1. OK.\\
Buscar chave inexistente & Ausente. OK.\\
Inserir \texttt{hash} novamente & Chave duplicada. OK.\\
Inserir \texttt{tabela} em colisão & Índices 0 e 0, uma colisão. OK.\\
Inserir terceira chave em m=2 & Tabela cheia. OK.\\
Remover \texttt{hash}, buscar \texttt{tabela} & \texttt{tabela} continua presente. OK.\\
Atualizar \texttt{tabela=7} & Mesma chave, frequência 7 e n=1. OK.\\\bottomrule
\end{tabular}
\medskip
\normalsize
\textbf{Oito casos conferidos pelo menu do hashapp.}\par
Registro em \arquivo{resultados/testes_manuais.md}: entradas, esperado, obtido e situação. A opção 6 exibe a estrutura. Redimensionamento não se aplica a E1.1.
\end{frame}

\begin{frame}{Análise experimental: metodologia}
\begin{enumerate}
\item Usar $n=100$, $1.000$ e $10.000$ registros, com $m=2n$ e $\alpha=0{,}5$.
\item Criar uma tabela nova em cada uma das \textbf{cinco repetições} por cenário.
\item Inserir as primeiras n chaves: \texttt{palavra0}, \texttt{palavra1}, etc.
\item Buscar as n chaves presentes e depois outras n chaves ausentes.
\item Medir inserção e buscas com \texttt{clock()}, além das colisões de inserção.
\end{enumerate}
\begin{block}{Regra de medição}
Tempos de CPU em milissegundos para o lote inteiro. A média é a soma das cinco medições dividida por cinco. Geração das chaves e criação da tabela ficam fora da medição.
\end{block}
\end{frame}

\begin{frame}{Resultados: médias de cinco repetições}
\begin{center}
\small\setlength{\tabcolsep}{7pt}\renewcommand{\arraystretch}{1.6}
\begin{tabular}{rrrrrrr}\toprule
n & m & $\alpha$ & Colisões & Inserção & Busca ok & Busca aus.\\
 & & & & (ms) & (ms) & (ms)\\\midrule
100 & 200 & 0,50 & 18 & 0,023 & 0,012 & 0,036\\
1.000 & 2.000 & 0,50 & 304 & 0,354 & 0,144 & 0,281\\
10.000 & 20.000 & 0,50 & 4.953 & 3,606 & 1,595 & 3,302\\
\bottomrule\end{tabular}
\end{center}
\begin{itemize}
\item No maior cenário, 4.953 inserções encontraram o índice inicial ocupado.
\item Capacidade suficiente: nenhuma necessidade de crescimento.
\item Tempos representam lotes de n operações. Para uma operação, dividir por n.
\item Os tempos podem variar quando os experimentos forem executados em outra máquina.
\end{itemize}
\end{frame}

\begin{frame}{Discussão dos resultados e limitações}
\begin{itemize}
\item Mesmo com $\alpha=0{,}5$, as colisões mudam com a capacidade e o conjunto de chaves.
\item A dobra pode concentrar índices. A sondagem linear pode formar blocos de posições ocupadas.
\item Uma busca ausente precisa encontrar um vazio ou percorrer a tabela, enquanto uma busca presente para ao encontrar a palavra.
\item As chaves são sintéticas. Os tempos não representam a leitura e normalização de um texto completo.
\item Mantivemos alfa constante: não medimos uma curva ao variar a ocupação.
\end{itemize}
\begin{alertblock}{Precisão da medição}
Lotes pequenos têm precisão limitada pela resolução de \texttt{clock()}. Repetimos o mesmo conjunto para comparar as medições de forma consistente.
\end{alertblock}
\end{frame}

\begin{frame}[fragile]{Demonstração 1: compilar com GCC}
Na pasta \texttt{tabela-hash}, executar:
\begin{lstlisting}[language=bash,basicstyle=\ttfamily\small]
gcc -std=c11 -Wall -Wextra -O2 \
  src/tabela.c src/aplicacao.c src/main.c -o hashapp

gcc -std=c11 -Wall -Wextra -O2 \
  src/tabela.c src/experimentos.c -o experimentos
\end{lstlisting}
\begin{itemize}
\item O primeiro comando gera \texttt{hashapp}, o programa interativo.
\item O segundo gera \texttt{experimentos}, usado somente para medir desempenho.
\item A barra no final da linha continua o comando no terminal.
\end{itemize}

\end{frame}

\begin{frame}[fragile]{Demonstração 2: colisão, duplicata e tabela cheia}
\begin{lstlisting}[language=bash,basicstyle=\ttfamily\small]
./hashapp 2
\end{lstlisting}
\small
\renewcommand{\arraystretch}{1.25}
\begin{tabular}{@{}p{1.8cm}p{5.5cm}p{5.8cm}@{}}\toprule
Opção & Respostas seguintes & Resultado esperado\\\midrule
1 & Palavra \texttt{hash}, frequência 1 & OK.\\
2 & Palavra \texttt{hash} & Frequência 1.\\
1 & Palavra \texttt{hash}, frequência 9 & Chave duplicada.\\
1 & Palavra \texttt{tabela}, frequência 2 & OK, com colisão.\\
6 & Nenhuma & m=2, n=2, uma colisão.\\
1 & Palavra \texttt{extra}, frequência 1 & Tabela cheia.\\\bottomrule
\end{tabular}
\medskip
\normalsize
No campo \textbf{Opção}, digitar o número. A palavra só deve ser digitada quando aparecer \textbf{Palavra}.
\end{frame}

\begin{frame}{Demonstração 3: remoção, atualização e listagem}
Continuar na mesma execução, após o exemplo anterior.
\medskip
\small\renewcommand{\arraystretch}{1.35}
\begin{tabular}{@{}p{1.3cm}p{5.2cm}p{6.6cm}@{}}\toprule
Opção & Respostas seguintes & Resultado esperado\\\midrule
4 & Palavra \texttt{hash} & OK.\\
2 & Palavra \texttt{tabela} & Presente, frequência 2.\\
2 & Palavra \texttt{hash} & Chave não encontrada.\\
6 & Nenhuma & Posição 0: REMOVIDO. Posição 1: tabela.\\
3 & Palavra \texttt{tabela}, frequência 7 & OK, mantém a chave.\\
5 & Nenhuma & Uma palavra: tabela, frequência 7.\\
0 & Nenhuma & Encerra o programa.\\\bottomrule
\end{tabular}
\medskip
\normalsize
\textbf{Linhas vazias são esperadas:} estes registros foram inseridos manualmente, sem arquivo de origem.
\end{frame}

\begin{frame}[fragile]{Demonstração 4: indexar o texto de exemplo}
\begin{lstlisting}[language=bash,basicstyle=\ttfamily\small]
./hashapp 64
\end{lstlisting}
\begin{enumerate}
\item Escolher a opção \textbf{7}.
\item Informar \texttt{dados/texto\_exemplo.txt}.
\item Conferir: \textbf{55 ocorrências, 41 palavras distintas, 0 rejeitadas}.
\item Escolher \textbf{2} e consultar \texttt{Hash}: chave \texttt{hash}, frequência \textbf{3}, linhas \textbf{1, 2 e 5}.
\item Escolher \textbf{5} para listar e \textbf{6} para exibir a estrutura.
\item Escolher \textbf{0} para sair.
\end{enumerate}
\begin{block}{O que este exemplo demonstra}
Normalização, repetição de palavras, frequência, linhas de ocorrência e consulta por chave.
\end{block}
\end{frame}

\begin{frame}[fragile]{Demonstração 5: executar os roteiros prontos}
\begin{lstlisting}[language=bash,basicstyle=\ttfamily\small]
./hashapp 64 < dados/demo_config40.txt

./hashapp 32 < dados/demo_tabela_cheia.txt
\end{lstlisting}
\begin{itemize}
\item O símbolo \texttt{<} usa as linhas do arquivo como respostas do menu.
\item Primeiro roteiro: indexação, duplicata, consultas, remoção, atualização e listagem.
\item Segundo roteiro: capacidade 32 para 41 palavras distintas.
\item Resultado do segundo: \textbf{44 ocorrências aceitas, 32 palavras distintas e 11 ocorrências rejeitadas}.
\end{itemize}
\begin{alertblock}{Ocorrências não são palavras distintas}
As 11 rejeitadas incluem repetições de palavras que não conseguiram entrar. Palavras já armazenadas continuam acumulando frequência.
\end{alertblock}
\end{frame}

\begin{frame}[fragile]{Demonstração 6: medições e registro manual}
\begin{lstlisting}[language=bash,basicstyle=\ttfamily\small]
./experimentos

./experimentos > resultados/experimentos_config40.md
\end{lstlisting}
\begin{itemize}
\item Correção: repetir os casos dos slides anteriores pelo menu.
\item Conferir entradas, esperado e obtido em \arquivo{resultados/testes_manuais.md}.
\item Experimentos: conferir três linhas de resultados, com médias de cinco repetições.
\item O símbolo \texttt{>} salva a saída, substituindo o conteúdo anterior do arquivo.
\item Novas execuções podem produzir tempos diferentes, mesmo com as mesmas chaves.
\end{itemize}
\end{frame}

\end{document}